\section{秦：关中还没有成为帝国}

后来的人知道秦会统一天下，于是很容易在春秋秦身上提前看见帝国的影子。可在平王东迁的年代，秦首先只是西部的一支力量：要面对戎狄，要经营关中，还要争取被中原诸侯承认。它离会盟中心很远，离边防、河谷和开垦却很近。

\eventanchor{SA-0770-QIN-ENFEOFF}{春秋·秦获封关中}

平王东迁时，秦襄公出兵护送。传统材料称周王室许秦有岐以西之地，条件是秦能自己收复。这样的封赏很有西周遗风：名分由王室给出，土地与秩序却要受命者用兵力和经营去取得。秦得到的不是一份已经划清边界的遗产，而是一项漫长的任务。关中有河谷、平原和向东的门户，也有需要长期应付的西部压力；秦人在此后的岁月里修城、用兵、与西戎周旋，才慢慢把“被封的名分”变成“能支配的腹地”。

\eventanchor{SA-0766-QIN-WEN-ASCEND}{春秋·秦文公继位}约在前766年，秦文公承接襄公留下的政治资本；约前763年，传世《史记·秦本纪》又把他在汧渭河谷卜居营邑写成经营西部的起点，\eventanchor{SA-0763-QIN-QIAN}{春秋·秦文公汧渭卜居}。这类早期年次和卜居叙事主要来自后出的国别本纪，与《竹书纪年》异文及早期都邑的历史地理研究相互参照，不能据此钉定城址或一次筑城的规模。不过，它保存了一个不宜略去的结构：护王的军政集团若不能安置在能供给、能通行的河谷节点，东迁所带来的名分便无从转成持续的资源。

\phantomsection\label{fig:vol01:spring-autumn:qin-jin-xiao-corridor}
\begin{center}
\begin{tikzpicture}[x=.88cm,y=.88cm,font=\small]
  \fill[paper] (0,0) rectangle (11.3,5.8);

  \draw[source!65!timeline,line width=1pt]
    (.35,4.85) .. controls (2.7,4.45) and (4.3,5.25) .. (6.7,4.7)
    .. controls (8.4,4.3) and (9.6,4.75) .. (10.9,4.35);
  \node[text=source!80!ink,above] at (5.8,4.8) {黄河方向};

  \fill[dynasty!19] (.75,1.35) -- (3.15,1.2) -- (3.45,3.15) -- (1.05,3.45) -- cycle;
  \node[align=center] at (2.15,2.35) {秦\\关中方向};
  \fill[timeline!19] (4.45,3.55) -- (6.65,3.5) -- (6.95,4.65) -- (4.9,4.9) -- cycle;
  \node[align=center] at (5.7,4.2) {晋\\河东方向};
  \fill[source!16] (4.25,1.45) -- (6.65,1.25) -- (7.2,2.8) -- (4.75,3.05) -- cycle;
  \node[align=center] at (5.7,2.1) {崤函山地\\通道方向};
  \fill[timeline!15] (8.2,1.0) -- (10.55,1.15) -- (10.8,3.15) -- (8.5,3.35) -- cycle;
  \node[align=center] at (9.55,2.15) {郑\\中原枢纽};

  \draw[-{Stealth[length=2mm]},ultra thick,color=dynasty]
    (3.35,2.55) .. controls (5.35,1.55) and (7.2,1.75) .. (8.05,2.05);
  \draw[-{Stealth[length=2mm]},thick,color=dynasty]
    (6.75,4.0) .. controls (7.75,3.7) and (8.05,3.05) .. (8.25,2.65);
  \node[text=dynasty,below,align=center] at (5.65,1.25) {前630年：秦、晋围郑\\暂时合军};
  \node[text=dynasty,right,align=left] at (7.25,3.55) {晋军\\南下};

  \draw[-{Stealth[length=2mm]},thick,dashed,color=timeline]
    (3.35,2.1) .. controls (4.6,2.65) and (6.55,2.7) .. (8.05,2.35);
  \node[text=timeline,above,align=center] at (5.7,2.95) {秦军东向、撤经崤一带\\（路径不可精确复原）};

  \draw[-{Stealth[length=2mm]},ultra thick,color=source]
    (5.9,3.5) -- (6.05,2.95);
  \fill[dynasty] (6.05,2.9) circle (.11);
  \node[text=source,above right,align=left] at (6.05,2.95)
    {前627年殽之战\\晋军阻击方向};

  \node[text=source,align=center,font=\footnotesize] at (2.45,.65)
    {山地通道使远征的补给、撤退\\与侧击风险彼此相连};
\end{tikzpicture}

\smallskip
\small\textit{图\quad 秦、晋、郑与崤函通道示意：将围郑、秦军东征郑与殽之战置于同一交通瓶颈中，帮助理解同盟裂缝怎样转为秦晋对抗。参考《春秋·僖公三十、三十三年》、《左传·僖公三十、三十二、三十三年》及《史记·秦本纪》《郑世家》；会战部署、具体行程和部分故事细节主要由后出叙事展开，尚有争议。所有边界与路线均为约略示意。}
\end{center}


东迁给秦打开的是机会，不是通行证。关中通向东面的河渡、山口和谷道，既让秦可接近诸侯政治，也使一支离开本土的军队要把粮食、车马与归路一道带进计算。传世材料并没有留下关中赋役、粮仓或道路的逐年簿籍，不能把这种约束改写成精确的运输图；但援立、输粟、围郑与殽道受挫连续留下的行动痕迹，足以提示资源与路线不是战事的背景，而是决定选择能否兑现的条件。

秦向东时，首先遇到的不是一片可以直取的中原，而是晋。两国隔着黄河与山道相望，既能以婚姻、援立和互市维持关系，也会因土地、粮食与通路彼此牵制。齐桓公死后，宋襄公曾试图以会盟填补中原的协调空缺；这只是理解当时诸侯秩序松动的一角。秦、晋之间更迫切的事情，是相邻国家能否在饥年、继承与边境争执中兑现盟友的承诺。

约在前650年，秦曾助流亡在外的夷吾返晋即位，\eventanchor{SA-0650-QIN-YIWU}{春秋·秦助夷吾返晋}。传世叙述把这次援立同河西土地的许诺相连，也把它写成日后反目的开端；较可把握的是，晋的继承危机给了秦介入东岸政治的机会，而援立不会消除双方对资源和边境的不同计算。秦以兵力帮助一个新政权站稳，所得到的并非可以永远收取的报偿，而是一段必须不断检验的关系。

\begin{event}{公元前647年}{秦向饥荒中的晋输送粟米}{经文年次可校，数量与言辞有争议}{秦、晋}
\eventanchor{SA-0647-QIN-JIN-GRAIN}{春秋·秦输粟援晋}

饥荒把这种关系放到最具体的尺度上：粟米能否运到，谁来承担运输与分配的成本，受援者又能否在下一次危机中回报。前647年，秦向歉收中的晋输粟。粮食不是会盟辞令，而是能够接续军政、贵族网络和乡里生计的物资；对晋惠公的政权而言，得到这批粮食能暂缓眼前的压力，对秦而言，救援也把关中安全和盟约信用系在晋国的稳定上。

次年形势倒转，秦也遭遇困难，晋却未能按传世材料所说的那样报偿。这里不宜把“输粟”写成一位君主的单纯仁慈，也不宜把未报偿简化为一场私人背信。现存材料容许我们看到的，是两个以农业产出、贵族征调和边境安全维持自身的国家，在歉收时发现盟约没有替它们准备共同的仓储或强制履约的办法。粮食暂时压住了危机，也把债务、互信和可供出兵的资源推到更尖锐的位置。
\end{event}

前645年，争执终于以兵车而非粮车收束于韩原。\eventanchor{SA-0645-HANYUAN}{春秋·韩原之战}中，秦军俘获晋惠公。秦并未因此得到一个可以任意支配的晋国；俘虏和议和首先说明，先前以援立与输粟维持的关系已经不能约束双方的军事选择。对秦而言，胜利显示它能够越过关中的门户影响东边；对晋而言，失败则把新政权在继承、边境和盟友信用上的脆弱一并暴露出来。

\begin{textwindow}[粮食、战报与后出的完整故事]
《春秋·僖公十五年》以“晋侯及秦伯战于韩，获晋侯”钉住韩原之战的结果，这是传世经文所提供的简短年代骨架。关于前647年输粟、次年未报偿以及援立条件的连贯经过，主要见于《左传·僖公十三年》《国语·晋语》及《史记·秦本纪》《晋世家》。这些后出的叙事保存了理解危机的线索，却也加入了粮食数量、政治言辞和人物评价；它们不能替代一份已经失传的仓储簿或外交盟书。故而可以确认救荒与战争相接，不能据此复原每一车粮食的路线，或断定每一方在当时只怀有一种目的。
\end{textwindow}

韩原之后的秦晋并未立刻成为两条绝不相交的线。前630年，晋、秦又共同围郑，\eventanchor{SA-0630-QINJIN-ZHENG}{春秋·秦晋围郑}。晋方面把郑在重耳流亡时的拒绝视为可用的旧账；秦方面则要衡量，攻下一座东向要城是否真比保留一个能沟通中原的郑国更合算。围城使两国短暂站在同一营垒，却没有抹去各自的利益。秦随后退兵，郑得以保存，联合行动中已经显出裂缝。

《春秋·僖公三十年》保留了“晋人、秦人围郑”的简笔记事。烛之武入说秦伯的完整言辞，则出自《左传》的叙事安排，后来又成为最著名的春秋外交场景之一。它准确地抓住了秦在郑问题上的利益困境，却不能当作夜间谈判的逐字记录。秦为何退师，既可从郑的中介价值与秦的东进盘算解释，也可能还受围城成本、晋国力量和当时信息限制影响；现存材料不足以让我们把这些因素排出唯一的先后。

晋文公死后，秦没有放弃向东的道路。前626年，秦穆公遣军东征郑，\eventanchor{SA-0626-QIN-ATTACK-ZHENG}{春秋·秦东征郑}。这支远征军要穿过或绕过晋国能够施加影响的地带，补给线因而同攻城本身一样危险。郑仍是中原诸侯关系的枢纽，绕过晋去经营郑，或许能重新打开秦的东向影响；但从现存材料看，秦军究竟以速取郑城、恢复盟约，还是以更广的东进布局为主，不能断言。把后世关于郑商人与预警的故事当作确定的情报流程，只会遮住远征本已显见的风险。

\begin{event}{公元前627年}{秦军在殽道受晋军击败}{可靠纪年}{殽}
\eventanchor{SA-0627-XIAO}{春秋·殽之战}

秦军未能把东征转成对郑的持久控制，归途中在殽山一带遭晋军与姜戎击败。《春秋·僖公三十三年》所记“晋人及姜戎败秦师于殽”，给出战败者、参战者和战场名称；它没有交代一支军队在狭窄通道中每一步如何失去主动。可以确定的是，山道把距离、补给与侧击的危险集中到一起，晋则能把对本土通路的熟悉和盟友资源用于阻截。

殽之战不是秦“从此不能东进”的判决。它更直接地表明，依靠一次远征越过晋的势力范围，尚不足以建立可持续的中原位置；秦晋间由救荒、援立和围城反复检验过的关系，至此转成更公开的对抗。晋守住了东向通道的优势，秦也必须重新估量把有限人力和粮食投往何处。
\end{event}

\begin{debate}[殽道的确定与东征的限度]
《左传·僖公三十二、三十三年》与《史记·秦本纪》把东征、退兵和战败组织为一段因果鲜明的故事，是理解此役的主要叙事材料。殽山及相关道路的大体战略意义没有疑问，但具体战场位置、行军路线和伏击经过仍须依赖历史地理考订；秦军是否已经形成足以稳定占领郑地的计划，也不能由后出故事单独证明。所谓“秦东进受挫”，说的是这次行动的结果与通道风险，不是把秦此后所有东向政策都预写成失败。
\end{debate}

殽后不久，晋把反制推进到秦的东缘。前625年，\eventanchor{SA-0625-PENGYA}{春秋·晋率诸侯伐秦取彭衙}晋率宋、陈、郑等国伐秦，取汪、彭衙而还；经文可校其年次，攻取地点的精确位置与战果规模则仍有争议。联军能抵达关中东部，却没有把一次攻取变成对秦的持久支配。次年，\eventanchor{SA-0624-QIN-WANGGUAN}{春秋·秦伐晋取王官}秦又渡河伐晋，取王官、郊等地后还师。关于焚舟等富于修辞的行动细节，不宜拿来绘制确切的行军路线；较为清楚的是，黄河渡口与沿线城邑控制使双方都能越界施压，却很难把短促远征固定为稳定占有。

东路受牵制后，秦把更多力量放回关中西部。前623年前后，\eventanchor{SA-0623-QIN-WEST}{春秋·秦西向扩张}所记的西向经营，显示秦并未因殽之败退出诸侯竞争，而是在另一面继续取得土地、人口和边防资源。后世“霸西戎”的说法容易把很长的接触、战争和兼并压缩成一位君主的功业；西部诸群体的称谓、范围和被控制的程度都仍有争议。可以较稳妥地说，东向道路受晋牵制，是秦加强西部经营的重要条件之一，并非唯一原因。

\eventanchor{SA-0621-QIN-MU-DEATH}{春秋·秦穆公卒，康公即位}前621年，穆公去世，康公接过的并不是一条已经安定的“西进国策”，而是仍与晋相接、又须处置西部关系的关中政局。前615年的\eventanchor{SA-0615-HEQU}{春秋·河曲之战}中，晋军又击败秦师；经文可确立这一胜负，战场位置和将帅评价仍主要依赖《左传》叙事。河曲失利使晋暂保河东与中原的联络线，也让秦晋竞争更像长期的边境消耗，而非一次决战的结局。把前面的援立、救荒、围郑、殽道和西向经营一概归结为“穆公称霸”，是后出的整齐读法；《左传》《国语》与《史记》分别保存的，是年次、言说传统和回望中的国别长线。它们能让人看见一位君主在这些选择中的位置，却不足以复原他的全盘谋划或继承安排。

\eventanchor{SA-0608-QIN-GONG-ASCEND}{春秋·秦康公卒，共公继位}约在前608年，康公死、共公继位的年次主要由后出的本纪和公系整理维系；边境军队、城邑和朝廷权力如何交接，并没有留下可逐项核对的记录。能够提出的只是继承带来的实际问题：新君须接手既有的河西压力，而不是从一张空白的外交图上选择方向。次年，\eventanchor{SA-0607-QIN-JIAO}{春秋·秦伐晋围焦}秦伐晋并围焦，赵盾救焦后又联宋、卫、陈侵郑；《春秋》留下交战的年次，而围焦、救焦与郑宋战事的完整行军顺序主要见于叙事材料，不宜由此断言诸国已有一套协调周密的战线。它至少显出，晋须在西线救援与中原盟友之间调兵，多方向战争不断拉扯霸主可用的资源。前580年，晋、秦会于令狐，\eventanchor{SA-0580-LINGHU}{春秋·晋秦会令狐}并恢复婚姻关系；具体婚姻安排与盟约文本已不可全知，但既有联姻网络至少提供了战后重新接触的政治语言。

秦晋的边境压力也没有随穆公而止。前578年，\eventanchor{SA-0578-MASUI}{春秋·麻隧之战}晋军败秦；《春秋》与《左传》可使这次交战进入年次骨架，战术和参战细节则主要属于叙事材料。此役不宜写成秦从此被逐出东方，却同殽之战一样提醒人们：通向中原的道路既是机会，也是让对手能够集中阻截的走廊。战场地望与行军路径仍须由历史地理研究考订，不能由一则完整战事故事倒推出唯一地点。

此后的经文把西线的紧张保留成几笔短而硬的记号：\eventanchor{SA-0563-JIN-ATTACK-QIN}{春秋·晋师伐秦}在前563年，\eventanchor{SA-0562-QIN-ATTACK-JIN}{春秋·秦人伐晋}在前562年，前559年又有\eventanchor{SA-0559-JIN-ATTACK-QIN}{春秋·晋士匄率诸侯伐秦}。它们不提供攻伐范围、伤亡或入秦深度，不能拼接成一场可以逐日复原的全面战争；但反复用兵本身说明，晋在郑、宋、楚方向经营联盟时，仍不得不为西部边界调度道路与粮秣，秦也始终在这一牵制中寻找余地。

秦并非只从晋的方向理解自己的位置。吴军在\eventref{SA-0506-BOJU}{春秋·柏举之战}攻入郢都后，前505年秦出兵援楚，\eventanchor{SA-0505-QIN-CHU-AID}{春秋·秦援楚收复郢都}，助楚逐步恢复核心都城与政权连续性。申包胥“哭秦庭”的完整场景见于《左传》《国语》及后出的叙事传统，不能据以还原求援与决策的每一步；较可把握的是，楚的腹地并未随郢都失守而消失，秦的援助与吴军难以久驻陌生区域共同限制了胜利的转化。前501年\eventanchor{SA-0501-QIN-AI-DEATH}{春秋·秦哀公去世}，经文又留下秦君交接的节点；继位礼与朝廷权力如何分配无从细究，不能因为后来的秦更强便为这一时刻预设答案。

\begin{turningpoint}[制度得失：盟约无法替代粮道与通道]
秦晋关系先后以援立、输粟、俘虏、围城和伏击来检验。援助解决了晋国一时的粮食与政权危机，代价是把报偿、边境和安全问题留在一份没有强制执行者的盟约里；联合围郑一度降低了彼此直接冲突的成本，却又让郑的价值和东进的路线显露为分歧。殽之败迫使秦承认，关中的国家能力不能只靠一次东征证明。它转而经营西部，在边防、土地和人口上积累资源，同时仍须应付晋的反复进攻与中原局势的牵引；这些积累日后能被战国改革接入新的制度，并不意味着春秋的秦已经预知统一的道路。
\end{turningpoint}

\begin{sourcenote}[秦早期的叙事陷阱]
《春秋》经文为韩原、围郑、麻隧及后来的秦晋攻伐保留了可校验的简短年次；《左传》《国语》展开了危机中的人物、言辞与选择；《史记·秦本纪》则在秦统一又灭亡之后，给予早期秦一条格外连贯的上升线。三者互读，能够看见粮食、通道和联盟如何相互牵动；不能把后出的完整叙事倒灌为现场实录，更不能把秦穆公、五羖大夫等富于文学力量的故事当作统一事业的预告。殽、麻隧及其道路的地望考订，同样只能收窄行动的空间，不能把伏击、行军或君主意图补成唯一版本。写春秋秦，应先看到它如何在西部经营、向东受挫，再谈它后来如何进入天下。
\end{sourcenote}
